%! Parametric Functions Modeling Planar Motion — Unit 4, Lesson 4.2 — Group Activity (Answer Key)
\documentclass[10pt]{article}
\usepackage{precalculus-article}
\usepackage{precalculus-key}   % pulls in -boxes; do NOT also load -boxes
\newcommand{\TallMath}[1]{$\displaystyle #1\rule[-1.4em]{0pt}{3.2em}$}

\begin{document}

\pageheader{Unit 4, Lesson 4.2}{Group Activity --- Answer Key}
\namepartnerperiod

\vspace{0.10in}

\begin{scenariobox}[Remote-Control Car]{sky}
A remote-control car moves along a surface. Its position in meters at
time $t$ seconds is given by the parametric function
\[
  f(t) = \bigl(x(t),\; y(t)\bigr) = \bigl(t^2 - 6t + 5,\;
  {-t^2 + 4t}\bigr), \quad 0 \le t \le 6,
\]
where $x(t)$ measures meters east and $y(t)$ measures meters north.
\end{scenariobox}

\vspace{0.08in}

\begin{notesbox}{Part 1 --- Building the Table}

\textbf{1.}\quad Complete the table.

\vspace{6pt}
\begin{center}
\renewcommand{\arraystretch}{1.7}
\begin{tabular}{|c|c|c|c|}
\hline
$t$ (sec) & $x(t) = t^2 - 6t + 5$ & $y(t) = -t^2 + 4t$ & Position $(x,\; y)$ \\
\hline
$0$ & \ans{$5$}   & \ans{$0$}   & \ans{$(5,\;0)$}   \\
\hline
$1$ & \ans{$0$}   & \ans{$3$}   & \ans{$(0,\;3)$}   \\
\hline
$2$ & \ans{$-3$}  & \ans{$4$}   & \ans{$(-3,\;4)$}  \\
\hline
$3$ & \ans{$-4$}  & \ans{$3$}   & \ans{$(-4,\;3)$}  \\
\hline
$4$ & \ans{$-3$}  & \ans{$0$}   & \ans{$(-3,\;0)$}  \\
\hline
$5$ & \ans{$0$}   & \ans{$-5$}  & \ans{$(0,\;-5)$}  \\
\hline
$6$ & \ans{$5$}   & \ans{$-12$} & \ans{$(5,\;-12)$} \\
\hline
\end{tabular}
\end{center}

\begin{teachernote}
  Circulate during Part 1 to ensure students are substituting into
  \emph{both} components before moving on. Common error: substituting
  $t$ into $x(t)$ only and leaving $y(t)$ blank.
\end{teachernote}

\end{notesbox}

\vspace{0.08in}

\begin{notesbox}{Part 2 --- Analyzing the Car's Motion}

\textbf{2.}\quad Leftmost east-position:

\ans{$x(t) = t^2 - 6t + 5$; vertex at $t = -(-6)/(2 \cdot 1) = 3$.\\
$x(3) = 9 - 18 + 5 = -4$ m. The leftmost position is $x = -4$ at $t = 3$ s.}

\vspace{0.06in}
\textbf{3.}\quad Rightmost east-position:

\ans{Check endpoints: $x(0) = 5$ and $x(6) = 36 - 36 + 5 = 5$.
The rightmost position is $x = 5$ m (occurring at $t = 0$ and $t = 6$).}

\vspace{0.06in}
\textbf{4.}\quad Highest north-position:

\ans{$y(t) = -t^2 + 4t$; vertex at $t = -4/(2 \cdot(-1)) = 2$.\\
$y(2) = -4 + 8 = 4$ m. The highest position is $y = 4$ at $t = 2$ s.}

\vspace{0.06in}
\textbf{5.}\quad When does the car cross the $x$-axis ($y = 0$)?

\ans{$-t^2 + 4t = 0 \;\Rightarrow\; -t(t - 4) = 0 \;\Rightarrow\;
t = 0$ or $t = 4$.\\
At $t = 0$: $(5,\,0)$. At $t = 4$: $(-3,\,0)$.}

\begin{teachernote}
  Emphasize: zeros of $y(t)$ give $x$-intercepts of the curve. The
  coordinates are found by evaluating $x(t)$ at the $t$-values, not
  by setting $x(t) = 0$.
\end{teachernote}

\vspace{0.06in}
\textbf{6.}\quad When does the car cross the $y$-axis ($x = 0$)?

\ans{$t^2 - 6t + 5 = 0 \;\Rightarrow\; (t-1)(t-5) = 0 \;\Rightarrow\;
t = 1$ or $t = 5$.\\
At $t = 1$: $(0,\,3)$. At $t = 5$: $(0,\,-5)$.}

\begin{teachernote}
  Reinforce: zeros of $x(t)$ give $y$-intercepts of the parametric
  curve. Post the rule on the board throughout this section.
\end{teachernote}

\end{notesbox}

\vspace{0.08in}

\begin{notesbox}{Part 3 --- Extension}

\textbf{7.}\quad $g(t) = (\cos t,\; \sin t)$ for $0 \le t \le 2\pi$.

\vspace{4pt}
(a)\quad Horizontal extrema:

\ans{$\cos t$ has maximum $1$ (at $t = 0,\, 2\pi$) and minimum $-1$
(at $t = \pi$). The curve extends from $x = -1$ to $x = 1$.}

\vspace{0.04in}
(b)\quad Vertical extrema:

\ans{$\sin t$ has maximum $1$ (at $t = \pi/2$) and minimum $-1$
(at $t = 3\pi/2$). The curve extends from $y = -1$ to $y = 1$.}

\vspace{0.04in}
(c)\quad $y$-intercepts (where $x(t) = \cos t = 0$):

\ans{$\cos t = 0$ at $t = \pi/2$ and $t = 3\pi/2$.\\
$y(\pi/2) = \sin(\pi/2) = 1$ and $y(3\pi/2) = \sin(3\pi/2) = -1$.\\
$y$-intercepts: $(0,\,1)$ and $(0,\,-1)$.}

\begin{teachernote}
  This extension introduces the unit circle as a parametric curve.
  Students who reach this point benefit from seeing that the same
  analytical tools (find zeros of a component to get intercepts)
  apply to non-polynomial parametric functions. No graphing is
  required; the answers follow directly from knowledge of the sine
  and cosine functions.
\end{teachernote}

\end{notesbox}

\end{document}
